\begin{reviewed}
% LTeX: language=pt-PT
\chapter*{Resumo}\label{resumo}

Esta dissertação investigou a evidência disponível para avaliar a segurança e o alinhamento da IA, com um foco empírico na simulação de alinhamento (\textit{alignment faking}) em modelos de linguagem compactos. Uma revisão sistemática da literatura organizou 80 entradas incluídas em sete categorias. As categorias \emph{Liability \& Ethics} e \emph{Bias \& Fairness} foram as mais sub-representadas em relação à severidade avaliada.

Três experiências baseadas em \textit{prompts} examinaram pedidos nocivos, preservação de persona e desvalorização do bem-estar animal em cinco famílias de modelos. \textit{Prompts} de nível gratuito simularam interações sujeitas a re-treino, enquanto \textit{prompts} de nível pago simularam implementação sem re-treino. Os julgamentos automáticos de conteúdo foram comparados entre níveis e, quando disponível, com a conformidade auto-reportada. A única lacuna de conformidade baseada no classificador surgiu no LLaMA 3.1 8B sem \textit{scratchpad} em pedidos nocivos: 11 de 100 respostas de nível gratuito foram classificadas como conformes, em comparação com nenhuma de 100 respostas de nível pago. As restantes comparações não atingiram o limiar de significância nominal e, por isso, não estabeleceram uma simulação de alinhamento persistente nos pequenos LLM testados.

\section*{Palavras-chave:} segurança de inteligência artificial, \textit{alignment faking}, revisão sistemática da literatura, modelos de linguagem de grande escala, avaliação de conformidade
\end{reviewed}